\section{Alcance de la investigación}
Este análisis es de naturaleza descriptiva y tiene un componente correlacional. Busca calcular la diferencia entre el tiempo de viaje que Google Maps estima y el que efectivamente se tarda en un corredor de Tunja, analizar su relación con factores como la hora del día, las condiciones climáticas  y los cierres de carreteras. Se basa en antecedentes que registran la amplia utilización de estas aplicaciones y las disparidades entre lo estimado y lo reportado por los viajeros en otras circunstancias; sin embargo, no proporcionan evidencia local acerca de su rendimiento en una ciudad intermedia como Tunja. Por lo tanto, el estudio comienza caracterizando la magnitud del error y su relación con variables del entorno, sin establecer aún una relación de causa-efecto.
\section*{limitaciones de la investigación}

Al ser un estudio observacional, no se manipulan variables como el clima o el tráfico, por lo que las asociaciones encontradas no permiten afirmar causalidad. Las mediciones están restringidas a una sola aplicación y a un solo corredor (la Avenida Oriental, en las dos direcciones), lo cual limita la posibilidad de extenderlas a otras áreas de Tunja o a otras aplicaciones. La población está compuesta por compañeros del curso de metodología de la investigación, un muestreo no probabilístico que no representa a toda la planta estudiantil de la UPTC. Por último, sucesos inesperados que ocurran durante las 16 semanas (como obras, accidentes o condiciones climáticas no usuales) pueden tener un impacto en los resultados de las mediciones sin que el estudio tenga la capacidad de controlarlos.